% ════════════════════════════════════════════════════════════════
\subsection{Problem sumy podzbiorów}
% ════════════════════════════════════════════════════════════════

\begin{problem}[sumy podzbiorów (Subset-Sum)] \label{problem:subsetsum}
Instancja: $S = \{x_1, x_2, \ldots, x_n\} \subseteq \N$, $t \in \N$.
\begin{itemize}[nosep]
	\item \emph{Wersja decyzyjna:} czy istnieje $S' \subseteq S$ takie, że $\sum_{x_i \in S'} x_i = t$?
	\item \emph{Wersja optymalizacyjna:} znaleźć $S' \subseteq S$ takie, że $\sum_{x_i \in S'} x_i$ jest największą możliwą liczbą $\leq t$.
\end{itemize}
\end{problem}

Dla zbioru $P \subseteq \N$ i $x \in \N$ oznaczamy $P + x = \{s + x : s \in P\}$.\\
np.\ dla $P = \{1, 2, 3, 6\}$ i $x = 2$: \quad $P + x = \{3, 4, 5, 8\}$.

Dla listy $L$ liczb naturalnych i $x \in \N$ przez $L + x$ oznaczamy listę powstałą z $L$ przez dodanie $x$ do każdego elementu.\\
np.\ $L = (3, 5, 8)$ i $x = 3$: \quad $L + x = (6, 8, 11)$.

\begin{lemma} \label{lemma:subsetsum-Pi}
	Niech $S = \{x_1, \ldots, x_n\} \subseteq \N$. Niech $P_0 = \{0\}$ i $P_i = P_{i-1} \cup (P_{i-1} + x_i)$ dla $i \in [n]$. Zbiór $P_n$ składa się ze wszystkich możliwych sum podzbiorów zbioru $S$.
\end{lemma}
\begin{proof}
	Pokażemy przez indukcję po $i$, że zbiór $P_i$ składa się ze wszystkich możliwych sum elementów ze zbioru $\{x_1, x_2, \ldots, x_i\}$.

	$i = 0$: \quad $P_0 = \{0\}$ i $\sum_{x_i \in \emptyset} x_i = 0$. $\checkmark$

	Załóżmy, że $P_{i-1}$ składa się ze wszystkich możliwych sum elementów z $\{x_1, \ldots, x_{i-1}\}$.\\
	Niech $S' \subseteq \{x_1, \ldots, x_i\}$ będzie dowolnym podzbiorem o~sumie $\sigma = \sum_{x_k \in S'} x_k$. Pokażemy, że $\sigma \in P_i$.
	\begin{itemize}[nosep]
		\item Jeśli $x_i \notin S'$, to $S' \subseteq \{x_1, \ldots, x_{i-1}\}$, więc $\sigma = \sum_{x_k \in S'} x_k \in P_{i-1} \subseteq P_i$.
		\item Jeśli $x_i \in S'$, to $S' \setminus \{x_i\} \subseteq \{x_1, \ldots, x_{i-1}\}$, więc
		      \begin{flalign*}
			      \sigma = \sum_{x_k \in S'} x_k = \sum_{x_k \in S' \setminus \{x_i\}} x_k + x_i \in P_{i-1} + x_i \subseteq P_i &  &  &  &
		      \end{flalign*}
	\end{itemize}
	Zatem $P_i$ zawiera wszystkie możliwe sumy elementów z $\{x_1, \ldots, x_i\}$.\\
	$P_i$ nie zawiera innych elementów, bo $P_{i-1}$ składa się z sum elementów z $\{x_1, \ldots, x_{i-1}\}$ a $P_{i-1} + x_i$ zawiera tylko sumy elementów z $\{x_1, \ldots, x_{i-1}, x_i\} \subseteq \{x_1, \ldots, x_i\}$.\\
	Dla $i = n$ mamy tezę.
\end{proof}

Lemat \ref{lemma:subsetsum-Pi} podsuwa algorytm: wystarczy liczyć kolejne zbiory $P_i$. Reprezentujemy je posortowanymi listami $L_i$ bez duplikatów, a rekurencję $P_i = P_{i-1} \cup (P_{i-1} + x_i)$ realizujemy scalaniem list. Dodatkowo z $L_i$ usuwamy elementy $> t$ -- żadna suma większa od $t$ nie może być rozwiązaniem, więc nie wpływa to na wynik.

Załóżmy, że mamy procedurę $\Call{MergeLists}{L, L'}$, która dla posortowanych list liczb naturalnych $L$, $L'$ zwraca posortowaną listę elementów z $L$ i~$L'$ i~usuwa duplikaty. Złożoność czasowa $\Call{MergeLists}{L, L'}$ wynosi $\BT(|L| + |L'|)$.

\begin{alg}{Exact Subset Sum}{alg:exactsubsetsum}
	\FunctionNN{ExactSubsetSum}{$S, t$}
	\State $n \Assign |S|$
	\State $L_0 \Assign (0)$
	\For{$i \Assign 1$ \To $n$}
	\State $L_i \Assign \Call{MergeLists}{L_{i-1}, L_{i-1} + x_i}$
	\State usuń z $L_i$ wszystkie liczby $> t$
	\EndFor
	\State \Return największy element $L_n$
	\EndFunction
\end{alg}

Poprawność: indukcyjnie $L_i$ to posortowana lista tych elementów $P_i$, które są $\leq t$. Na mocy lematu \ref{lemma:subsetsum-Pi} $P_n$ zawiera wszystkie sumy podzbiorów $S$, więc największy element $L_n$ jest największą taką sumą $\leq t$ -- czyli rozwiązaniem wersji optymalizacyjnej.

\paragraph{Złożoność czasowa}
Długość listy $L_i$ może być wykładnicza względem $i$, nawet $2^i$. Dla $I = (S, t)$:
\begin{flalign*}
	|I| = \BT\left(\log t + \sum_{i=1}^{n} \log x_i\right) &  &  &  &
\end{flalign*}
Jeśli liczby $t$ i $x_i$ są $n$-cyfrowe, to $\log t, \log x_i = \BT(n)$ i~wtedy $|I| = \BT(n + n^2) = \BT(n^2)$. Skoro $L_n$ może mieć długość $2^n$, to \Call{ExactSubsetSum}{} (Algorytm~\ref{alg:exactsubsetsum}) jest algorytmem wykładniczym.

Jeśli liczba $t$ jest ograniczona przez funkcję wielomianową od $n = |S|$, czyli $t = \BO(n^p)$ dla jakiegoś $p \in \N$, to długość $L_i$ jest ograniczona przez $\BO(t) = \BO(n^p)$ i w tym przypadku algorytm działa w czasie wielomianowym.

\paragraph{Schemat aproksymacyjny}
Pokażemy w~pełni wielomianowy schemat aproksymacyjny dla problemu sumy podzbiorów (Problem~\ref{problem:subsetsum}).

\paragraph{Pomysł:} będziemy skracać listy $L_i$. Dla $0 < \delta < 1$ z listy $L$ robimy krótszą listę $L'$ taką, że dla każdego elementu $y$ usuniętego z $L$ istnieje element $z$ w $L'$, który jest blisko niego, dokładnie $\frac{y}{1 + \delta} \leq z \leq y$.

\begin{example}[Procedura Trim]
	$L = (15, 16, 17, 25, 27, 28, 35, 36, 37)$, $\delta = 0{,}1$ (czyli próg $1 + \delta = 1{,}1$).

	Idziemy po liście od lewej; ostatni zachowany element nazwijmy $z$. Kolejny $y$ zachowujemy tylko wtedy, gdy $y > z \cdot 1{,}1$ -- w~przeciwnym razie $y$ jest na tyle blisko $z$, że $z$ go reprezentuje ($\frac{y}{1{,}1} \leq z \leq y$).
	\begin{itemize}[nosep]
		\item $15$ -- pierwszy element, zachowujemy; $z = 15$.
		\item $16 \leq 15 \cdot 1{,}1 = 16{,}5$ -- usuwamy, reprezentuje je $15$.
		\item $17 > 16{,}5$ -- zachowujemy; $z = 17$.
		\item $25 > 17 \cdot 1{,}1 = 18{,}7$ -- zachowujemy; $z = 25$.
		\item $27 \leq 25 \cdot 1{,}1 = 27{,}5$ -- usuwamy, reprezentuje je $25$.
		\item $28 > 27{,}5$ -- zachowujemy; $z = 28$.
		\item $35 > 28 \cdot 1{,}1 = 30{,}8$ -- zachowujemy; $z = 35$.
		\item $36 \leq 35 \cdot 1{,}1 = 38{,}5$ -- usuwamy, reprezentuje je $35$.
		\item $37 \leq 38{,}5$ -- usuwamy, reprezentuje je $35$.
	\end{itemize}
	Stąd $L' = (15, 17, 25, 28, 35)$.
\end{example}

Następująca procedura $\Call{Trim}{L, \delta}$ realizuje ten pomysł.\\
$L = (y_1, \ldots, y_m)$ -- posortowana rosnąco lista, $0 < \delta < 1$.

\begin{alg}{Trim}{alg:trim}
	\FunctionNN{Trim}{$L, \delta$}
	\State $m \Assign |L|$
	\State $L' \Assign (y_1)$
	\State $last \Assign y_1$
	\For{$i \Assign 2$ \To $m$}
	\If{$y_i > last \cdot (1 + \delta)$}
	\State dołącz $y_i$ na koniec $L'$
	\State $last \Assign y_i$
	\EndIf
	\EndFor
	\State \Return $L'$
	\EndFunction
\end{alg}

Czas działania $\Call{Trim}{L, \delta}$ wynosi $\BT(|L|)$.

Oto schemat aproksymacyjny:

\begin{alg}{Approx Subset Sum}{alg:approxsubsetsum}
	\FunctionNN{ApproxSubsetSum}{$S, t, \varepsilon$}
	\State $n \Assign |S|$
	\State $L_0 \Assign (0)$
	\For{$i \Assign 1$ \To $n$}
	\State $L_i \Assign \Call{MergeLists}{L_{i-1}, L_{i-1} + x_i}$
	\State usuń z $L_i$ liczby $> t$
	\State $L_i \Assign \Call{Trim}{L_i, \frac{\varepsilon}{2n}}$ \Comment{$\delta = \frac{\varepsilon}{2n}$}
	\EndFor
	\State \Return największy element $z \in L_n$ \Comment{$z^*$}
	\EndFunction
\end{alg}

\begin{example}[ApproxSubsetSum od początku do końca]
	$S = \{104, 102, 201, 101\}$ (w~tej kolejności $x_1, \ldots, x_4$), $t = 308$, $\varepsilon = 0{,}4$. Wtedy $n = 4$ i~próg przycinania to $\delta = \frac{\varepsilon}{2n} = \frac{0{,}4}{8} = 0{,}05$ (czyli mnożnik $1 + \delta = 1{,}05$). Startujemy z~$L_0 = (0)$.

	W~każdej iteracji najpierw scalamy $L_{i-1}$ z~$L_{i-1} + x_i$, potem usuwamy elementy $> t = 308$, a~na końcu przycinamy. Przy przycinaniu zachowujemy element $y$, gdy $y > last \cdot 1{,}05$, gdzie $last$ to ostatnio zachowany element.
	\begin{itemize}[nosep]
		\item \textbf{$i = 1$, $x_1 = 104$.} \\
		      scalenie: $(0) \cup (104) = (0, 104)$; \\
          po usunięciu $> 308$: $(0, 104)$; \\
          przycinanie: \\
		      \begin{tblr}{colspec={r c l l}, column{4}={font=\itshape}, rowsep=1pt}
			      $0$   &       &                       & zachowane, $last = 0$   \\
			      $104$ & $>$   & $0 \cdot 1{,}05 = 0$   & zachowane, $last = 104$ \\
		      \end{tblr} \\
		      $L_1 = (0, 104)$.
		\item \textbf{$i = 2$, $x_2 = 102$.} \\
		      scalenie: $(0, 104) \cup (102, 206) = (0, 102, 104, 206)$;\\
          po usunięciu $> 308$: bez zmian; \\
          przycinanie: \\
		      \begin{tblr}{colspec={r c l l}, column{4}={font=\itshape}, rowsep=1pt}
			      $0$   & $ $       &                          & zachowane, $last = 0$   \\
			      $102$ & $>$      & $0$                      & zachowane, $last = 102$ \\
			      $104$ & $\leq$   & $102 \cdot 1{,}05 = 107{,}1$ & usunięte             \\
			      $206$ & $>$      & $107{,}1$                & zachowane, $last = 206$ \\
		      \end{tblr} \\
		      $L_2 = (0, 102, 206)$.
		\item \textbf{$i = 3$, $x_3 = 201$.} \\
		      scalenie: $(0, 102, 206) \cup (201, 303, 407) = (0, 102, 201, 206, 303, 407)$; \\
          po usunięciu $> 308$: $(0, 102, 201, 206, 303)$; \\
          przycinanie: \\
		      \begin{tblr}{colspec={r c l l}, column{4}={font=\itshape}, rowsep=1pt}
			      $0$   & $ $       &                              & zachowane, $last = 0$   \\
			      $102$ & $>$      & $0$                          & zachowane, $last = 102$ \\
			      $201$ & $>$      & $102 \cdot 1{,}05 = 107{,}1$ & zachowane, $last = 201$ \\
			      $206$ & $\leq$   & $201 \cdot 1{,}05 = 211{,}05$ & usunięte               \\
			      $303$ & $>$      & $211{,}05$                   & zachowane, $last = 303$ \\
		      \end{tblr} \\
		      $L_3 = (0, 102, 201, 303)$.
		\item \textbf{$i = 4$, $x_4 = 101$.} \\
		      scalenie: $(0, 102, 201, 303) \cup (101, 203, 302, 404) = (0, 101, 102, 201, 203, 302, 303, 404)$; \\
          po usunięciu $> 308$: $(0, 101, 102, 201, 203, 302, 303)$; \\
          przycinanie: \\ 
		      \begin{tblr}{colspec={r c l l}, column{4}={font=\itshape}, rowsep=1pt}
			      $0$   & $ $       &                              & zachowane, $last = 0$   \\
			      $101$ & $>$      & $0$                          & zachowane, $last = 101$ \\
			      $102$ & $\leq$   & $101 \cdot 1{,}05 = 106{,}05$ & usunięte               \\
			      $201$ & $>$      & $106{,}05$                   & zachowane, $last = 201$ \\
			      $203$ & $\leq$   & $201 \cdot 1{,}05 = 211{,}05$ & usunięte               \\
			      $302$ & $>$      & $211{,}05$                   & zachowane, $last = 302$ \\
			      $303$ & $\leq$   & $302 \cdot 1{,}05 = 317{,}1$ & usunięte                \\
		      \end{tblr} \\
		      $L_4 = (0, 101, 201, 302)$.
	\end{itemize}
	Zwracamy największy element $z = 302$ (podzbiór $\{201, 101\}$). Optimum to $z^* = 307$ (podzbiór $\{104, 102, 101\}$), wycięte przez przycinanie w~iteracji $i = 4$. Iloraz $\frac{z^*}{z} = \frac{307}{302} \approx 1{,}017 \leq 1 + \varepsilon = 1{,}4$ -- zgodnie z~gwarancją.
\end{example}

% ════════════════════════════════════════════════════════════════
\subsection{NP-zupełność problemu sumy podzbiorów}
% ════════════════════════════════════════════════════════════════

\begin{theorem}
	Problem sumy podzbiorów (Problem~\ref{problem:subsetsum}) jest NP-zupełny.
\end{theorem}
\begin{proof}
	Problem jest w NP. Dla instancji $(S, t)$ i $S' \subseteq S$ można sprawdzić w czasie wielomianowym czy $\sum_{x_i \in S'} x_i = t$.

	Pokażemy redukcję $\text{3-SAT} \leq_p \text{SUBSET-SUM}$.

	$\Phi$ -- formuła 3-SAT w CNF, $x_1, x_2, \ldots, x_n$ -- zmienne, $C_1, C_2, \ldots, C_m$ -- klauzule. Załóżmy, że żadna klauzula nie zawiera zmiennej i jej zaprzeczenia (wpp.\ jest spełniona) oraz każda zmienna występuje w co najmniej jednej klauzuli.

	\begin{dygresja}
		Przypomnijmy: formuła jest w~postaci normalnej koniunkcyjnej (CNF), gdy jest koniunkcją (iloczynem) klauzul $\Phi = C_1 \wedge \cdots \wedge C_m$, a~każda klauzula jest alternatywą (sumą) literałów, czyli zmiennych lub ich zaprzeczeń, np.\ $C_j = (x_a \vee \overline{x_b} \vee x_c)$. Dopisek ,,3-'' oznacza, że każda klauzula ma dokładnie trzy literały.
	\end{dygresja}

	Utworzymy instancję $(S, t)$ problemu SUBSET-SUM dla $\Phi$.\\
	Dla każdej zmiennej $x_i$ mamy w $S$ dwie liczby $r_i$ i $r_i'$, $(n+m)$-cyfrowe.\\
	Dla każdej klauzuli $C_j$ mamy w $S$ dwie liczby $s_j$ i $s_j'$, $(n+m)$-cyfrowe.\\
	$t$ jest liczbą $(n+m)$-cyfrową: $\underbrace{1\,1\,\ldots\,1}_{n}\,\underbrace{4\,4\,\ldots\,4}_{m}$.

	Każdą z tych liczb traktujemy jak ciąg $n + m$ cyfr: pierwsze $n$ pozycji etykietujemy zmiennymi $x_1, \ldots, x_n$, kolejne $m$ pozycji -- klauzulami $C_1, \ldots, C_m$. Cyfrę na danej pozycji ustawiamy tak:
	\begin{itemize}[nosep]
		\item \textbf{liczby zmiennych} $r_i$ (dla $x_i = 1$) oraz $r_i'$ (dla $x_i = 0$): obie mają $1$ na pozycji $x_i$, a~$0$ na pozostałych pozycjach zmiennych. Na pozycji klauzuli $C_j$ mają $1$, jeśli dany literał występuje w~$C_j$ ($x_i$ dla $r_i$, $\overline{x_i}$ dla $r_i'$), w~przeciwnym razie $0$.
		\item \textbf{liczby klauzul} $s_j$ i~$s_j'$: mają $0$ wszędzie poza pozycją $C_j$, gdzie $s_j$ ma $1$, a~$s_j'$ ma $2$. Służą jako ,,dopełnienie'' pozwalające dobić pozycję $C_j$ do żądanej cyfry $4$.
	\end{itemize}
	Schematycznie (każda komórka to cyfra danej liczby na danej pozycji):
	\begin{center}
		\begin{tblr}{colspec={l|cc|cc}, row{1}={font=\itshape}, vline{2,4}, hline{2}}
			       & pozycja $x_i$ & inne poz.\ zmiennych & pozycja $C_j$                    & inne poz.\ klauzul \\
			$r_i$  & $1$           & $0$                  & $1$ gdy $x_i \in C_j$            & $0$                \\
			$r_i'$ & $1$           & $0$                  & $1$ gdy $\overline{x_i} \in C_j$ & $0$                \\
			$s_j$  & $0$           & $0$                  & $1$                              & $0$                \\
			$s_j'$ & $0$           & $0$                  & $2$                              & $0$                \\
		\end{tblr}
	\end{center}
	Liczby $r_i, r_i', s_j, s_j'$ są parami różne -- z konstrukcji i założeń o $\Phi$.\\
	Na każdej pozycji suma wszystkich liczb $\leq 6$, więc dodawanie odbywa się cyfra po cyfrze, bez przeniesień.

	\begin{dygresja}
		Skąd cyfry $1$ i~$4$ w~$t$, i~dlaczego wymuszają one spełnialność? Bez przeniesień każdą pozycję rozpatrujemy osobno.

		\emph{Pozycje zmiennych (cel $1$).} Jedyne liczby z~$1$ na pozycji $x_i$ to $r_i$ i~$r_i'$. Aby suma na tej pozycji wyniosła dokładnie $1$, trzeba wziąć do $S'$ dokładnie jedną z~nich -- to właśnie jest wybór wartości $x_i$ ($r_i \leftrightarrow x_i = 1$, $r_i' \leftrightarrow x_i = 0$).

		\emph{Pozycje klauzul (cel $4$).} Na pozycji $C_j$ wybrane liczby zmiennych dają sumę $k \in \{0, 1, 2, 3\}$ -- to liczba literałów $C_j$ spełnionych przez wartościowanie. Brakującą resztę $4 - k$ możemy dobrać liczbami dopełniającymi $s_j$ ($=1$) i~$s_j'$ ($=2$), które dają do dyspozycji $0$, $1$, $2$ lub $3$ ($= 1 + 2$). Stąd:
		\begin{itemize}[nosep]
			\item $k = 1$: dobieramy $s_j$ i~$s_j'$ ($1 + 2 = 3$),
			\item $k = 2$: dobieramy $s_j'$ ($2$),
			\item $k = 3$: dobieramy $s_j$ ($1$).
		\end{itemize}
		Liczbę $4$ dobrano tak, by maksymalne dopełnienie ($3$) wystarczało wtedy i~tylko wtedy, gdy $k \geq 1$. Gdy $k = 0$ (klauzula niespełniona), brakuje $4$, a~dopełnienie sięga tylko $3$ -- pozycji nie da się uzupełnić, więc nie istnieje $S'$ o~sumie $t$.
	\end{dygresja}

	Redukcja działa w czasie wielomianowym: zbiór $S$ składa się z~$2(n + m)$ liczb $(n + m)$-cyfrowych, utworzenie każdej z~nich oraz liczby $t$ zajmuje czas liniowy względem $n + m$.

	$\Phi$ spełnialna $\implies$ istnieje spełniające wartościowanie.\\
	Jeśli $x_i = 1$, to do $S'$ bierzemy $r_i$.\\
	Jeśli $x_i = 0$, to do $S'$ bierzemy $r_i'$.\\
	$C_j$ spełniona $\implies$ suma $r_i, r_i' \in S'$ ma na pozycji $C_j$ $1$ lub $2$ lub $3$.\\
	Dokładamy do $S'$ odpowiednio $s_j$ i $s_j'$ lub $s_j'$ lub $s_j$.\\
	Wtedy $S'$ ma sumę równą $t$.

	$S' \subseteq S$ ma sumę równą $t$ $\implies$ do $S'$ należy dokładnie jedna z liczb $r_i, r_i'$:
	\begin{itemize}[nosep]
		\item jeśli $r_i \in S'$, to $x_i = 1$,
		\item jeśli $r_i' \in S'$, to $x_i = 0$.
	\end{itemize}
	$C_j$ jest spełniona, bo $s_j + s_j' \leq 3$ na pozycji $C_j$ $\implies$ co najmniej $1$ jest z $r_i, r_i'$ należących do $S'$ mająca $1$ na pozycji $C_j$:
	\begin{itemize}[nosep]
		\item $r_i \in S'$ to $x_i$ jest w $C_j$, a mamy $x_i = 1$,
		\item $r_i' \in S'$ to $\overline{x_i}$ jest w $C_j$, a mamy $x_i = 0$.
	\end{itemize}
\end{proof}

\begin{example}[Redukcja 3-SAT do SUBSET-SUM]
	$C_1 = (x_1 \vee \overline{x_2} \vee \overline{x_3})$, \quad $C_2 = (\overline{x_1} \vee \overline{x_2} \vee \overline{x_3})$,\\
	$C_3 = (\overline{x_1} \vee \overline{x_2} \vee x_3)$, \quad $C_4 = (x_1 \vee x_2 \vee x_3)$.\\
	$\Phi = C_1 \wedge C_2 \wedge C_3 \wedge C_4$. \quad $n = 3$, $m = 4$.

	Wszystkie liczby są $7$-cyfrowe: $3$ pozycje zmiennych i~$4$ pozycje klauzul. Cel:
	\begin{aztable}
		\centering
		\renewcommand{\arraystretch}{1.2}
		\begin{NiceTabular}{c|ccc|cccc}[margin,hvlines]
			    & $x_1$ & $x_2$ & $x_3$ & $C_1$ & $C_2$ & $C_3$ & $C_4$ \\
			\hline
			$t$ & 1     & 1     & 1     & 4     & 4     & 4     & 4     \\
		\end{NiceTabular}
	\end{aztable}

	\emph{Liczby zmiennych.} Para $r_i, r_i'$ ma $1$ na pozycji $x_i$ (więc do sumy $1$ trafia dokładnie jedna z~nich). Na pozycjach klauzul $1$ pojawia się tam, gdzie odpowiedni literał występuje w~klauzuli -- np.\ $r_1$ ($x_1$) ma $1$ w~$C_1$ i~$C_4$, bo $x_1 \in C_1, C_4$.
	\begin{aztable}
		\centering
		\renewcommand{\arraystretch}{1.2}
		\begin{NiceTabular}{c|ccc|cccc}[margin,hvlines]
			       & $x_1$ & $x_2$ & $x_3$ & $C_1$ & $C_2$ & $C_3$ & $C_4$ \\
			\hline
			$r_1$  & 1     & 0     & 0     & 1     & 0     & 0     & 1     \\
			$r_1'$ & 1     & 0     & 0     & 0     & 1     & 1     & 0     \\
			$r_2$  & 0     & 1     & 0     & 0     & 0     & 0     & 1     \\
			$r_2'$ & 0     & 1     & 0     & 1     & 1     & 1     & 0     \\
			$r_3$  & 0     & 0     & 1     & 0     & 0     & 1     & 1     \\
			$r_3'$ & 0     & 0     & 1     & 1     & 1     & 0     & 0     \\
		\end{NiceTabular}
	\end{aztable}

	\emph{Liczby dopełniające.} Para $s_j, s_j'$ ma niezerowe cyfry ($1$ i~$2$) tylko na pozycji $C_j$ -- służy do dobicia tej pozycji do~$4$.
	\begin{aztable}
		\centering
		\renewcommand{\arraystretch}{1.2}
		\begin{NiceTabular}{c|ccc|cccc}[margin,hvlines]
			       & $x_1$ & $x_2$ & $x_3$ & $C_1$ & $C_2$ & $C_3$ & $C_4$ \\
			\hline
			$s_1$  & 0     & 0     & 0     & 1     & 0     & 0     & 0     \\
			$s_1'$ & 0     & 0     & 0     & 2     & 0     & 0     & 0     \\
			$s_2$  & 0     & 0     & 0     & 0     & 1     & 0     & 0     \\
			$s_2'$ & 0     & 0     & 0     & 0     & 2     & 0     & 0     \\
			$s_3$  & 0     & 0     & 0     & 0     & 0     & 1     & 0     \\
			$s_3'$ & 0     & 0     & 0     & 0     & 0     & 2     & 0     \\
			$s_4$  & 0     & 0     & 0     & 0     & 0     & 0     & 1     \\
			$s_4'$ & 0     & 0     & 0     & 0     & 0     & 0     & 2     \\
		\end{NiceTabular}
	\end{aztable}

	\emph{Pełny obraz.} Weźmy spełniające wartościowanie $x_1 = 1$, $x_2 = 0$, $x_3 = 0$: bierzemy $r_1$ (bo $x_1 = 1$) oraz $r_2', r_3'$ (bo $x_2 = x_3 = 0$), a~na każdej pozycji klauzuli dobieramy $s_j, s_j'$ tak, by uzbierać $4$. Wybrane liczby ($S'$) zaznaczono; ostatni wiersz to ich suma.
	\begin{aztable}
		\centering
		\renewcommand{\arraystretch}{1.2}
		\begin{NiceTabular}{c|ccc|cccc}[margin,hvlines,color-inside]
			                          & $x_1$ & $x_2$ & $x_3$ & $C_1$ & $C_2$ & $C_3$ & $C_4$ \\
			\hline
			$t$                       & 1     & 1     & 1     & 4     & 4     & 4     & 4     \\
			\hline
			\rowcolor{gray!18} $r_1$  & 1     & 0     & 0     & 1     & 0     & 0     & 1     \\
			$r_1'$                    & 1     & 0     & 0     & 0     & 1     & 1     & 0     \\
			$r_2$                     & 0     & 1     & 0     & 0     & 0     & 0     & 1     \\
			\rowcolor{gray!18} $r_2'$ & 0     & 1     & 0     & 1     & 1     & 1     & 0     \\
			$r_3$                     & 0     & 0     & 1     & 0     & 0     & 1     & 1     \\
			\rowcolor{gray!18} $r_3'$ & 0     & 0     & 1     & 1     & 1     & 0     & 0     \\
			\rowcolor{gray!18} $s_1$  & 0     & 0     & 0     & 1     & 0     & 0     & 0     \\
			$s_1'$                    & 0     & 0     & 0     & 2     & 0     & 0     & 0     \\
			$s_2$                     & 0     & 0     & 0     & 0     & 1     & 0     & 0     \\
			\rowcolor{gray!18} $s_2'$ & 0     & 0     & 0     & 0     & 2     & 0     & 0     \\
			\rowcolor{gray!18} $s_3$  & 0     & 0     & 0     & 0     & 0     & 1     & 0     \\
			\rowcolor{gray!18} $s_3'$ & 0     & 0     & 0     & 0     & 0     & 2     & 0     \\
			\rowcolor{gray!18} $s_4$  & 0     & 0     & 0     & 0     & 0     & 0     & 1     \\
			\rowcolor{gray!18} $s_4'$ & 0     & 0     & 0     & 0     & 0     & 0     & 2     \\
			\hline
			$\sum$                    & 1     & 1     & 1     & 4     & 4     & 4     & 4     \\
		\end{NiceTabular}
	\end{aztable}
	Suma wybranych liczb równa się $t$, więc $S' = \{r_1, r_2', r_3', s_1, s_2', s_3, s_3', s_4, s_4'\}$ jest rozwiązaniem.
\end{example}
